\section{Generalized path integral and Riccati compression}\label{sec:riccati}

\begin{remark}[Generalized path integral and Riccati compression]
The acronym GPI in the in-context GPI theorem (\cref{thm:in-context-gpi}) admits a second reading:
\emph{Generalized Path Integral}.
The two readings coincide.

The forward estimator of \cref{prop:forward-only} supplies the
inner evaluation: for a given gait parameter~$\beta$, sample $M$
domain-randomized trajectories and average to obtain $J(\beta)$
(\cref{alg:riccati-eval}).
The Boltzmann target
$\pi(\beta) \propto \exp\!\bigl(-J(\beta)/\varepsilon\bigr)$
converts the uniform (prior) measure over trajectories to the
optimal (posterior) measure, with the cost functional $J$
playing the role of the negative log-likelihood.
The re-cycle outer loop (\cref{alg:recycle}) replaces the
Monte Carlo convergence rate $O(N^{-1/2})$ with the
information-theoretic rate $R_T/T \to 0$ via
$\gamma_T$ (\cref{prop:gp-regret}).

\medskip\noindent\textbf{Riccati compression.}
Specialize to $G = (\R^n, +)$ with linear dynamics
$\dot x = Ax + Bu$ and quadratic cost
$\int_0^T (x^\top Q\,x + u^\top R\,u)\,dt$
(the LQG regime),
where $A \in \R^{n \times n}$, $B \in \R^{n \times m}$,
$Q \succeq 0$, $R \succ 0$, and the pair $(A, B)$ is stabilizable
with $(A, Q^{1/2})$ detectable.
The Feynman--Kac path integral is Gaussian and evaluates in closed
form; the Boltzmann-weighted average compresses to the matrix Riccati
ODE:
\[
\dot P = A^\top\! P + PA - PBR^{-1}B^\top\! P + Q,
\]
where $P(t) \in \R^{n \times n}$ is the cost-to-go matrix and
$V(x) = x^\top P\, x$ the value function.
No sampling is needed: the exact solution replaces the Monte Carlo
estimator entirely.
The dual equation (Kalman--Bucy filter) compresses the estimation
filtration in the same way.
The CLT rate $O(N^{-1/2})$ of \cref{prop:forward-only} thus quantifies
the cost of leaving the Riccati regime: on a nonlinear Lie group~$G$
with non-quadratic cost, the path integral cannot be evaluated
analytically and one pays $O(N^{-1/2})$ per sample for the Monte Carlo
approximation.
\end{remark}

\medskip\noindent\textbf{FONC and SOSC: prior and posterior Riccati.}
The Boltzmann prior-to-posterior conversion above
(uniform measure $\to$ $\pi \propto e^{-J/\varepsilon}$)
has an exact counterpart in the backward pass.
Two variants of the Riccati recursion arise, distinguished by
their treatment of dynamics curvature---the second-order terms
$V'_x \cdot f_{xx}$ in the Q-function expansion.

\begin{remark}[FONC and SOSC backward passes]
\label{rem:fonc-sosc}
The DDP backward pass expands
$Q(x_t, u_t) = \ell_t + V_{t+1}(f_t)$ to second order.
The $Q_{uu}$ block decomposes as
\begin{equation}\label{eq:fonc-sosc}
Q_{uu}^{\mathrm{SOSC}}
= \underbrace{\ell_{uu} + f_u^\top V'_{xx}\, f_u}_{%
    Q_{uu}^{\mathrm{FONC}} \;\succeq\; 0}
  \;+\; \underbrace{\lambda^{*\top}_{t+1}\,\nabla^2_{uu} f_t}_{%
    \text{dynamics curvature}},
\end{equation}
where $V'_x = \lambda^*_{t+1}$ at optimality
(similarly for $Q_{xx}$ and $Q_{ux}$).
The \emph{first-order necessary conditions} (FONC) backward pass
drops the dynamics curvature, recovering
iLQR~\cite{TassaErezTodorov,MastalliCrocoddyl}:
$Q_{uu}^{\mathrm{FONC}} \succeq 0$ unconditionally---a
\emph{prior} on the cost landscape that ignores
how nonlinear the dynamics are.
The \emph{second-order sufficient conditions} (SOSC) backward pass
retains it, updating the prior with the dynamics Lagrangian Hessian
$\lambda^{*\top}\nabla^2 f$---a \emph{posterior} that is
positive-definite only on the dynamically reachable subspace
(the reduced Hessian
$Q_{uu}^{\mathrm{feas}}
= Q_{uu} - Q_{ux}\,Q_{xx}^{-1}\,Q_{xu} \succ 0$).
For \emph{linear} dynamics, $\nabla^2 f = 0$ and the prior
equals the posterior: FONC $=$ SOSC.
Frey et al.~\cite{FreyDiffMPC} prove the converse:
the QP (FONC) and NLP (SOSC) sensitivities
$\partial z^\star/\partial\theta$ coincide \emph{if and only if}
the Lagrangian Hessian is exact---equivalently, $\nabla^2 f = 0$.
For nonlinear dynamics, FONC sensitivities can be
``completely off''~\cite[Section~5.1]{FreyDiffMPC}.
\end{remark}

\begin{remark}[Implications for the re-cycle]
\label{rem:fonc-recycle}
The FONC/SOSC decomposition \eqref{eq:fonc-sosc} has three
consequences for the re-cycle architecture.
\begin{enumerate}[label=(\roman*),nosep]
\item \emph{Centroidal model.}
  The dynamics $\dot x = A(\sigma)\,x + B(\sigma)\,u + c$ is
  affine in $(x, u)$ for frozen contact~$\sigma$, so
  $\nabla^2 f = 0$ and FONC $=$ SOSC\@.
  The iLQR backward pass (\cref{alg:riccati-eval}) gives
  exact NLP sensitivities, not approximate ones---for
  both \cref{alg:quadruped} ($13$-state) and
  \cref{alg:action-rate} ($17$-state augmented).
\item \emph{GP-LCB architecture.}
  The re-cycle (\cref{alg:recycle}) evaluates $J(\theta)$
  as a black box: the GP posterior provides Bayesian
  ``sensitivities'' $(\mu_n, \sigma_n)$ through function
  evaluations, without computing
  $\partial z^\star/\partial\theta$ via the implicit function
  theorem.
  The two-level contraction (\cref{prop:unconditional})
  therefore holds unconditionally---it does not require
  the SOSC backward pass at either level.
\item \emph{Holonomy convergence.}
  The re-cycle converges to $H(f) = 0$
  (\cref{prop:student-convergence}), the Gaussian regime
  where the dynamics curvature correction vanishes and
  FONC $=$ SOSC\@.
  The distinction between prior and posterior Riccati
  dissolves at the fixed point.
\end{enumerate}
\end{remark}

